\documentclass[11pt,a4paper]{article}
\usepackage[T1]{fontenc}
\usepackage{lmodern}
\usepackage[margin=22mm]{geometry}
\usepackage{amsmath,amssymb,amsthm,microtype}
\usepackage[hidelinks]{hyperref}
\setlength{\parindent}{0pt}
\setlength{\parskip}{5pt}
\setlength{\emergencystretch}{2em}
\newtheorem{lemma}{Lemma}
\title{The Nine-Dimensional Snake Problem:\\A Short Explanation of the Current Argument}
\author{Iori Hironaka \quad Tommy Ma \quad Steven Zhan}
\date{Research state after round 220}
\begin{document}
\maketitle

\section{Question and evidence}
Let $a(n)$ be the largest edge length of an induced path in the binary
hypercube $Q_n$. Its vertices are distinct, and two are adjacent exactly
when their path indices are consecutive. A length-$m$ snake has $m+1$ vertices.

The maintained interval is
\[
\boxed{191\le a(9)\le233.}
\]
The lower bound comes from a verified construction. The contributor's
completed 233 theorem has a saved final assembly and standard-axiom audit,
with no remaining certificate or share-floor premise. Its source is pinned
to \texttt{b0521e1d} and audit to \texttt{b65d30ca}; a fresh independent
full rebuild remains pending~\cite{bound233archive}.

The known values $1,2,4,7,13,26,50,98$ do not determine the next term.
Numerical continuations suggest 194, 195 or 196, without proving a recurrence.
To establish 196 exactly requires both a 196-edge witness and exclusion
of every 197-edge snake; every longer snake contains such a prefix.
The following counting and construction arguments address different parts
of this task. Restricted search exclusions do not establish the global optimum.

\section{The square budget and the completed upper bound}
For a nontrivial snake put $v=m+1$, and let $s_j$ count coordinate squares
containing exactly $j$ selected vertices. Write $T=s_1$.
There are 4608 squares, each selected vertex lies in 36, and no square
is fully selected. Three-selected-corner squares correspond exactly to
internal path vertices, so $s_3=v-2$. Consequently
\[
s_0+s_1+s_2+s_3=4608,\qquad s_1+2s_2+3s_3=36v,
\]
and subtraction gives the exact identity
\[
\boxed{T+2s_0+35(m+1)=9214.}
\]

If $\alpha_i$ counts singleton squares at path vertex $i$, local geometry
gives positive counts, endpoint counts at least four, internal neighbours
of a count one at least four, and no three consecutive internal twos.
These restrictions imply sums at least 14 and 16 on six and seven
consecutive internal vertices. For example, every triple sums to at least
six, with equality only at $(1,4,1)$; the neighbouring count restriction
handles that exceptional triple when combining blocks.
The archived count-only potential and omitted boundary counts yield
$T\ge4m-13$, hence $39m+2s_0\le9192$. This historical route gives 235;
its potential retains a separate independent-rebuild requirement.

Empty squares retain information discarded by that count-only estimate.
For a cube vertex $q$, let $e(q)$ count its incident empty squares and
$d(q)$ its selected neighbours. Each selected vertex $x$ receives
\[
\sigma_{12}(x)=\sum_{q\sim x}
 \left\lfloor\frac{12e(q)}{d(q)}\right\rfloor.
\]
All denominators appearing here are positive. Let $L$ sum the remainders
$12e(q)\bmod d(q)$ for $d(q)>0$, and $12e(q)$ for $d(q)=0$.
Division with remainder and the four corners of every empty square give
\[
\sum_{x\in S}\sigma_{12}(x)+L=48s_0,\qquad
G+840(m+1)+L=221136,
\]
where $S$ is the selected vertex set and $G=24T+\sum_{x\in S}\sigma_{12}(x)$.

In the completed contributor comparison, its total window weight $W$
satisfies, for $m\ge7$,
\[
520(m-7)\le W+1023,\qquad W+4200m\le1099800.
\]
The checked local certificates, share floors and finite potential therefore give
\[
4720m\le1104463<4720\cdot234=1104480,
\]
so $m\le233$. Shorter paths already satisfy this bound. The saved assembly
reuses imported objects and cloud-built certificates; their fresh independent
rebuild remains distinct from the successful saved audit~\cite{bound233archive}.

The alternative retained comparison uses $m-8$ windows. Under its own
local-weight and state-walk hypotheses, boundary counting and telescoping give
\[
(4200+\mu)m+5L\le1098600+8\mu+R,
\]
where every state edge has weight plus potential drop at least $\mu$,
and the potential has range at most $R$. Its conditional 234 result,
and a separate relaxed 233 comparison, still require their stated local
certificates and geometric and graph interpretations. Neither improves the
maintained interval. Their unresolved premises do not belong to the completed
contributor theorem; detailed profiles remain in the companion manuscript.

\section{Joining lower-dimensional paths}
Cutting at the $k$ edges in one coordinate leaves $k+1$ path components
in the two faces, with total internal length $m-k$.
The saved representatives have decompositions
\[
26=2+13+9+2,\quad 50=18+24+6+2,\quad
98=8+25+10+25+26+4,
\]
where each final term counts crossings. These describe particular witnesses.

Paths of lengths $p,q$ in opposite faces join into a snake of length
$p+q+1$ if their projections meet only at their designated joining endpoint.
Cross-face adjacency requires identical projections, which proves the claim.
For a two-crossing construction, take disjoint, nonadjacent outer paths
$P,R$ in one face and a middle path $Q$ in the other. Each projected
intersection must consist of one endpoint of the corresponding outer path
and a different endpoint of $Q$. The resulting length is
$|E(P)|+|E(Q)|+|E(R)|+2$, explaining the checked $18+24+6+2=50$ example.

There is a restriction on modifying two positive-length parents $p,q$.
Delete $D$ original vertices and add $A$ new ones. Let $b$ count deleted
original endpoints, $e$ original edges deleted at both ends, and $H$ final
within-face edges touching a new vertex. If the result has length $m$
and $k$ crossings, vertex and edge counting give
\[
m+k+b+e+H=p+q+2+2A,\qquad m+k\le2a(n)+2+2A.
\]
Indeed, deletion removes $2D-b-e$ original edges, and the final vertex
count is $p+q+2-D+A=m+1$.
Without additions, a 196-edge lift from parents of length at most 98
uses lengths 98 and 97 without deletion, or two lengths 98 with one
deletion; it permits at most two crossings. These are necessary conditions.

The saved 98-edge $Q_8$ representative crosses every coordinate at least
four times. Its only four-crossing split has two 25-edge components in
one face, using 52 vertices, whereas a single $Q_7$ snake has at most 51.
The counting restriction excludes all other splits from a deletion-only
two-parent lift. This representative therefore already requires a more
general construction than deletion from two $Q_7$ parents.

\section{Changing the embedding of an entire snake}
Partition coordinates and write vertices as $(z,u)$. Assign an isometry
$g_z$ of the free-coordinate cube to each retained slice, and map
$(z,u)$ to $(z,g_z(u))$. Within-slice distances are preserved, and
nonadjacent slices cannot produce cube edges.
For adjacent slices with selected free-word sets $B_z,B_w$, put
$h=g_z^{-1}g_w$. All remaining edges and nonedges are preserved exactly when
\[
h(u)=u\quad(u\in B_z\cap B_w),\qquad
h(B_w\setminus B_z)\cap B_z=\varnothing.
\]
The conditions retain old crossings and forbid new projected coincidences.

This family produces a different cube class from the saved 50-edge parent.
For the saved 98-edge parent, the complete native census over all 254
nontrivial coordinate partitions produces no new class. Lean checks the
generic criterion and coverage of the literal finite tables. Their complete
geometric interpretation and the surviving leaves retain native status.

With two free coordinates, a written classification uses a uniform parity
offset followed by component switches; its complete Lean assembly is pending.
With three free coordinates, an injective partial map preserving parity,
adjacency and nonadjacency extends to an isometry: it permutes antipodal-pair
labels consistently. This extension fails in four coordinates, where equal
indexed parity patterns can have different endpoint distances. Arbitrary
coordinate histories therefore allow changes beyond slice isometries.

\section{Coordinate histories and finite capacity}
Retain the source base word $z_i$ and choose arbitrary free words $u_i$.
For Hamming distance $d$, the exact path conditions are
\[
d(z_i,z_j)+d(u_i,u_j)
\begin{cases}=1,&j=i+1,\\ \ge2,&j\ge i+2.\end{cases}
\]
They exclude both repeats and chords. Consecutive retained-coordinate steps
force a constant free word, so merge their maximal index intervals while
retaining every pair constraint. Each interval change flips one free bit;
interval word parities alternate. Translation, coordinate permutation and
lexicographic representatives remove symmetry without losing an orbit.

For this indexed 98-edge source, all 70 four-free-coordinate and all 56
five-free-coordinate models exhaust after source-isometric assignments are
excluded. This native result concerns the specified source and index matching.
For six free coordinates, retained pair $(1,4)$ also exhausts; 27 choices
remain unresolved. Four have saved unfinished searches, at $(0,4)$,
$(3,4)$, $(0,5)$ and $(0,6)$; the other 23 have only limited initial runs.
The latest $(0,6)$ continuation retains eight unfinished regions and no new parent.

With two retained coordinates, the base is a square. Distinct intervals of
the same free-word parity cannot have consecutive original indices.
If they share a word, every cross-interval base pair must be opposite.
A square corner has a unique opposite, so path injectivity forces both
intervals to be singletons at opposite corners. No three can share one word.

For one parity, let $M$ count intervals visiting multiple base vertices and
$N_z$ singleton intervals at corner $z$. Their distinct-word capacity gives
\[
\boxed{M+\max(N_{00},N_{11})+\max(N_{01},N_{10})\le2^{r-1}}
\]
for $r\ge1$ free coordinates: choose all multi-visit intervals and the larger
singleton family from each opposite pair. Every two chosen words differ.
For retained pair $(0,5)$ the even counts are $M=12$ and
$(N_{00},N_{01},N_{10},N_{11})=(5,4,10,10)$.
The chosen 32 intervals use every even word of $Q_6$; each of the nine
outside intervals must share with a chosen antipodal partner.

For a distinct-word family and a subset $J$, let $D_j$ be current domains
and $U=\bigcup_{j\in J}D_j$. Necessarily $|J|\le|U|$.
At equality, injectivity uses every word of $U$ within $J$, so other
members of that family must avoid $U$. Such removal requires the stated
distinctness and cannot automatically be applied across partner families.

The outside singleton family also has distinct words: its corners are equal
or adjacent. If $C$ is the chosen family and $O$ the outside family of one
parity in $Q_6$, at most $s=32-|C|$ outside values are absent from $C$.
For $(0,6)$, $s=2$ in even parity and $s=6$ in odd parity. Thus at least
seven of nine even and seven of thirteen odd outside intervals must share.
An outside interval with no domain intersection with any permitted partner
must consume one spare word. Once $s$ such intervals are forced, every
other outside interval must retain support in its partners' domains.

More generally, let $N(J)$ contain the chosen partners of outside intervals
$J$ whose domains intersect. Shared members inject into distinct partners,
so $|J|\le|N(J)|+s$. At equality, $J$ uses every value assigned to $N(J)$
and exhausts the spare allowance. Other outside intervals must therefore
share with chosen partners outside $N(J)$.

Lean checks the generic counting and assignment-preservation arguments.
For the literal $(0,6)$ source, checked index witnesses connect the actual
path to chosen/outside distinctness and complete partner lists. The normalized
initial word, explicit six-bit range and group boundaries also yield the
two 32-word parity pools. Current domain fitting remains an explicit premise;
native relation tables, matching discovery and complete traversal remain
separate from these checked path and filtering implications.

Consecutive triples supply another restriction. Two free steps leave the
outer words at distance zero or two, and each retained pair needs a compatible
third value in its current domain. For example, words 5 and 6 have only
common neighbours 4 and 7; removing both excludes that pair.
Lean checks preservation of every fitting assignment and validity under
smaller domains. The native continuation retains projected pair relations
within a job and reloads original relations for each new job.
No controlled comparison establishes a speed improvement.

Continuation itself relies on coverage: splitting $D$ into $\{a\}$ and
$D\setminus\{a\}$ partitions assignments. Saving unfinished descendants and
unexamined ancestor values preserves the unexplored region. Lean checks this
generic implication; saved-state audits do not certify the entire native search.

\section{A direct partner target and the fractional obstruction}
Let $P$ be the saved 98-edge $Q_8$ path and $t$ either endpoint, 0 or 78.
Remove $V(P)\setminus\{t\}$ and seek a 93-edge induced path beginning at $t$,
with its other endpoint free. The joining argument would give
$98+93+1=192$ edges in $Q_9$. This route does not require a different parent.
Each allowed graph has 158 vertices and 337 edges. Both direct searches
remain unresolved; the degree-and-cycle queries returned \texttt{UNKNOWN}.

A 93-edge partner starting at an even vertex has 47 vertices of each parity
and an odd final endpoint, as Lean derives from consecutive cube edges.
The model requires degree one at the two selected endpoints and degree two
at other selected vertices. It also uses the induced-edge relation
$y_{ab}=x_ax_b$ for Boolean vertex selections $x$.
Degree equations can admit separate cycles. An actual induced path must omit
a vertex of every cycle: its last cycle vertex cannot have two distinct
earlier cycle neighbours, since both would be its immediate predecessor.
Lean checks this argument and the resulting cycle occupancy inequality.

For any allowed vertex subset $U$ and $v\in U$, a stronger necessary inequality is
\[
\sum_{ab\in E(U)}y_{ab}\le\sum_{u\in U}x_u-x_v,
\]
where $E(U)$ consists of allowed edges with both endpoints in $U$.
The selected part of an actual path has fewer edges than vertices when
$v$ is selected; when it is unselected, the weaker inequality still holds.
Lean derives this from inducedness by assigning each selected edge to a
distinct selected path index, avoiding a selected marked index.
Writing $X,Y,B,E$ for selected vertices, internal edges, boundary edges
and selected endpoints in $U$, the degree sum gives
$2X-E=2Y+B$. Thus the same inequality is $B+E\ge2x_v$.

Relax the Boolean variables to $[0,1]$ and retain the linear product
inequalities, degree equations, cycle inequalities and subset inequalities.
Exact finite certificates give cardinality upper bounds 97.840921 at
root 0 and 97.359759 at root 78. Both are below 98, excluding a 97-edge
partner in this fixed-parent model; the conclusion is source-specific.
The old solver exclusions had no Lean contradiction certificate.

The new exact rational feasible cardinalities are
\[
q_0=\frac{22588258102073}{230867380070},\qquad
q_{78}=\frac{3545809907526}{36419686519}.
\]
Both exceed 94. To check all cycle inequalities, put $w_u=1-x_u$.
For each deleted edge $ab$, a rational potential has $p_a=0$,
$|p_u-p_v|\le w_u+w_v$ on other edges, and $p_b+w_a+w_b\ge2$.
Summing along the rest of any cycle gives $2\sum_{u\in C}w_u\ge2$,
which is the required occupancy bound. Separately, 158 feasible flows
per root certify every subset inequality: a flow of at least $2x_v$
from $v$ to an added sink is bounded by every cut separating them.
The cut capacity is exactly $B+E$. Mixing with the corresponding one-edge path
using weight $92/(q_t-2)$ gives total selection weight exactly 94.
All these linear inequalities are preserved by the mixture.
Thus even the full forest-and-cycle continuous model cannot exclude the
93-edge target. These fractional assignments do not construct a path.
The saved arrays have independent exact arithmetic audits; Lean checks the
generic degree, bound, cycle and forest arguments. Their numerical array
assembly remains separate~\cite{research220}.

Boolean consistency gives a further direction. On an allowed edge define
$\delta_{ab}=x_a+x_b-2y_{ab}$. For Boolean selections this records whether
the endpoint selections differ, so any cycle has an even number of such changes.
For a cycle $C$ and an odd subset $F$ of its edges, this implies
\[
\sum_{e\in F}\delta_e-\sum_{e\in E(C)\setminus F}\delta_e\le|F|-1.
\]
The saved root-0 point has disagreements $(0,1,0,0)$ around the square
$(44,45,61,60)$ and violates this inequality. These classical cycle-parity
constraints test consistency absent from the current relaxation~\cite{barahona1986}.
They have not yet been added to the linear model. Their effect on the best
bound, or on feasibility at 94 after reoptimization, remains unknown.

\section{Remaining questions}
A compatible 93-edge partner would improve the lower bound to 192.
A different parent or face arrangement remains possible, but finding a
new $Q_8$ parent alone does not supply a compatible $Q_9$ construction.
Exhausting all remaining six-coordinate models would establish only a
source-specific requirement of seven changed histories.
The maintained global interval remains $191\le a(9)\le233$, with an
independent rebuild of the completed upper-bound assembly pending.
Resolving the proposed value 196 still requires both its witness and
a universal 197-edge exclusion.

\smallskip
{\small Detailed proofs, witnesses and verification records are retained
in the companion manuscript and milestones through round 220.
The small-dimensional values appear in
\href{https://oeis.org/A099155}{OEIS A099155}; the published 191-edge
construction is in
\href{https://arxiv.org/html/2607.15270v3}{Orland et al., arXiv:2607.15270v3}.}

\begin{thebibliography}{9}
\bibitem{bound233archive}
Snake-in-the-box project contributors. \emph{Completed share-weighted bound
233: final assembly, standard-axiom audit and verification limitations}.
Contributor revision \texttt{b0521e1d14fd52c62db4cdeccc356b70a623eb82};
audit revision \texttt{b65d30ca15441420b8f4538f671e589c2502674f},
14 November 2026.
\url{https://github.com/hironaka-iori/Snake-in-the-box-in-9/tree/b0521e1d14fd52c62db4cdeccc356b70a623eb82/milestones/33-233-lean}.


\bibitem{research220}
I.~Hironaka, T.~Ma, and S.~Zhan. \emph{Subset forest bounds,
fractional feasibility and product consistency}. Research milestone 220,
15 November 2026.
\url{https://github.com/hironaka-iori/Snake-in-the-box-in-9/tree/main/milestones/220-subset-forest}.

\bibitem{barahona1986}
F.~Barahona and A.~R.~Mahjoub. On the cut polytope.
\emph{Mathematical Programming} \textbf{36} (1986), 157--173.
\url{https://www.lamsade.dauphine.fr/~mahjoub/Papers/BM86.pdf}.

\end{thebibliography}
\end{document}
